\documentclass[border=0pt]{standalone}
% The shared preamble of every docs diagram (docs/diagrams/<name>.tex).
% dvisvgm's driver first, so patterns and opacity survive into the SVG.
\def\pgfsysdriver{pgfsys-dvisvgm.def}
\usepackage[siunitx,RPvoltages]{circuitikz}
\usepackage{xcolor}
% Latin Modern: outline fonts for every glyph a label uses, µ included.
\usepackage[T1]{fontenc}
\usepackage{lmodern}
\usetikzlibrary{arrows.meta,backgrounds,calc,patterns,positioning}
\tikzset{
  % A white card, so diagrams read in mdBook's dark themes too.
  every picture/.append style={show background rectangle,
    background rectangle/.style={fill=white,inner sep=4pt}},
  % Layout geometry: layers told apart by fill and hatching, never colour alone.
  layer a/.style={draw=black,line width=0.6pt,fill=gray!20},
  layer b/.style={draw=black,line width=0.6pt,fill=gray!45},
  layer c/.style={draw=black,line width=0.6pt,fill=white,postaction={pattern=north east lines}},
  layer d/.style={draw=black,line width=0.6pt,fill=white,postaction={pattern=crosshatch dots}},
  marker/.style={draw=black,dashed,line width=0.6pt},
  % A measured distance and its label.
  dim/.style={{Latex[length=2mm]}-{Latex[length=2mm]},line width=0.5pt},
  dim label/.style={font=\small,fill=white,inner sep=1pt},
  % Where a rule fails: thick, dashed and labelled, not only red.
  violation/.style={draw=red!75!black,line width=1.4pt,dash pattern=on 3pt off 2pt},
  violation label/.style={font=\small\bfseries,text=red!75!black},
  caption/.style={font=\small},
}

\begin{document}
\begin{circuitikz}
  % Cross-section: two taps on a two-layer silicon profile.
  \fill[layer a] (0,0) rectangle (8,-0.9);
  \fill[layer b] (0,-0.9) rectangle (8,-2.6);
  \draw[line width=1.2pt] (0,-2.6) -- (8,-2.6);
  \node[caption,anchor=west] at (8.1,-0.45) {\SI{3}{\micro\metre}, \SI{10}{\ohm.cm}};
  \node[caption,anchor=west] at (8.1,-1.75) {\SI{300}{\micro\metre}, \SI{0.01}{\ohm.cm}};
  \node[caption,anchor=west] at (8.1,-2.6) {backside (ground)};
  \fill[layer c] (1.5,0) rectangle (2.5,0.3);
  \fill[layer c] (5.5,0) rectangle (6.5,0.3);
  \node[caption,above] at (2,0.3) {t1 (tap)};
  \node[caption,above] at (6,0.3) {t2 (tap)};
  \draw[dim] (2.5,0.15) -- (5.5,0.15) node[dim label,midway] {\SI{10}{\micro\metre}};
  % The port network GPurify extracts.
  \draw (1,-4) node[circ,label=above:{t1}]{} to[R, l=\SI{1.78}{\mega\ohm}] (7,-4) node[circ,label=above:{t2}]{};
  \draw (1,-4) to[R, l_=\SI{47.3}{\kilo\ohm}] (1,-5.8) node[ground]{};
  \draw (7,-4) to[R, l=\SI{47.3}{\kilo\ohm}] (7,-5.8) node[ground]{};
\end{circuitikz}
\end{document}
